%!TEX root = ../../thesis.tex
\chapter{先行研究}
\label{chap:prior_work}

本章では，本研究の基礎となる，本研究室のシナリオに基づくナビゲーションの研究について述べる．

\section{トポロジカルマップとシナリオ}
\label{sec:prior_topological}

% TODO: 島田らのトポロジカルマップとシナリオ（人の道案内の調査，通路の特徴と行動の組）を書く．

\section{シナリオ駆動型ナビゲーション}
\label{sec:related_scenario}

シナリオ駆動型ナビゲーション\cite{scenario_si2023}は，人間が道案内に用いるような文でロボットに経路を与える手法である．
シナリオは「突き当りまで直進．左折．」のように，通路の特徴と行動を組み合わせた文の列として記述する．
ロボットはシナリオを条件と行動の組に分解し，カメラ画像から認識した通路形状が条件を満たしたときに次の行動へ移る．
これにより，目的地までの経路を，距離や座標ではなく通路の形状の並びとして与えることができる．

従来手法では，通路形状の認識と，通路に沿った走行に，視覚と行動のend-to-end学習によるモデルを用いていた\cite{scenario_si2023}．
これらのモデルは走行対象の環境で収集したデータで学習するため，外観の異なる環境で用いるには，その環境でのデータ収集と再学習が必要であった．

\section{先行研究の課題}
\label{sec:prior_issue}

% TODO: 従来手法（春山）で用いたモデルの構成と，未学習環境での性能の課題を具体的に書く．
